% =====================================================
% Literature Review
% =====================================================
\section{Literature Review}
% TODO: Provide an overarching narrative of related work in DR grading.

\subsection{Model Selection}
% TODO: Discuss prior architectures (Inception-v3, ResNet, EfficientNet, DORN, etc.)
% and justify the chosen backbone.

\subsection{Preprocessing}
Preprocessing is not a cosmetic step in medical imaging---it is the single largest lever between a model that learns lesion-level pathology and one that latches onto background artefacts. Raw fundus photographs vary enormously in illumination, contrast, white-balance, the size of the visible retinal disc, and the amount of black border surrounding it; without standardisation, the network is forced to spend representational capacity on these nuisance factors instead of on disease cues~\cite{salehi2023preprocessing,litjens2017survey}. A consistent finding in the medical CNN literature is that, in the absence of preprocessing, gradient-based attribution maps drift onto clinically irrelevant regions (background, large vessels, image borders), whereas appropriate preprocessing concentrates attention on the diagnostic structures themselves~\cite{salehi2023preprocessing}.

Within our pipeline, preprocessing is built around four mutually reinforcing components. \textbf{Geometric normalisation} resizes every image to a fixed $224\times224$ resolution and (where applied) crops the circular field of view, removing the black corners that otherwise contribute strong but uninformative gradients. \textbf{Photometric normalisation} subtracts the per-channel ImageNet mean $(0.485,0.456,0.406)$ and divides by its standard deviation $(0.229,0.224,0.225)$, ensuring that the feature statistics fed to the pretrained ResNet match those it was trained on---without this step, transfer learning is markedly less effective. \textbf{Data augmentation} (horizontal/vertical flips, small rotations, brightness--contrast jitter, optional CLAHE contrast enhancement) synthetically expands the effective training distribution and is one of the most powerful regularisers known for limited medical datasets~\cite{shorten2019aug,perez2017aug}. \textbf{Regularisation hyperparameters} that complement preprocessing---most notably \emph{dropout}, which randomly disables units during training to prevent co-adaptation~\cite{srivastava2014dropout}, and \emph{weight decay}, which penalises large parameter norms---directly address the overfitting we observed in the baseline. Closely related are the choice of \emph{weight initialisation} (we rely on the He/Kaiming scheme that is the default for ReLU networks~\cite{he2015delving}), the optimiser learning rate, and the use of pretrained ImageNet weights. Together, this stack converts a brittle classifier that memorises the majority class into one whose validation curve tracks the training curve, and whose Grad-CAM attributions concentrate on hemorrhages, microaneurysms, and exudates rather than on the dataset's incidental statistics.

\subsection{Loss Function Definition}
% TODO: Compare cross-entropy, weighted CE, focal loss, ordinal regression losses,
% and CORAL framework.

\subsection{Evaluation Metric}
% TODO: Justify Quadratic Weighted Kappa as the primary evaluation metric;
% discuss accuracy, macro-F1, and confusion-matrix-based analyses.
